\chapter{Lesson 2: Tone Sandhi and the Writing System}

\section{Tone sandhi}

\emph{Tone sandhi} (\zh{变调} biàndiào) means that some syllables change tone because of
the tone that follows. Textbooks and dictionaries write the \emph{original} tone; we are
expected to pronounce the changed one. Only three cases matter at this level.

\subsection{Third tone + third tone}

\begin{important}[Rule 1: 3 + 3 $\rightarrow$ 2 + 3]
When two third tones come together, the first is pronounced as a second tone.
\begin{itemize}
  \item \zh{你好} written \emph{nǐ hǎo}, said \emph{ní hǎo}.
  \item \zh{很好} written \emph{hěn hǎo}, said \emph{hén hǎo}.
  \item \zh{洗澡} \emph{xǐzǎo} ``to bathe'' is said \emph{xízǎo}; \zh{可以} \emph{kěyǐ}
        ``may'' is said \emph{kéyǐ}.
\end{itemize}
\end{important}

Before any other tone (1st, 2nd, 4th or neutral) a third tone becomes a \emph{half third
tone}: it only dips low and does not rise. Example: \zh{我们} \emph{wǒmen}, \zh{老师}
\emph{lǎoshī}, \zh{好吗} \emph{hǎo ma}.

With three third tones in a row, the result depends on how the words group. In
\zh{我很好} \emph{wǒ hěn hǎo}, ``\zh{很好}'' is a unit, so we usually say
\emph{wǒ hén hǎo}; in fast speech \emph{wó hén hǎo} is also common.

\subsection{\zh{不} bù}

\begin{important}[Rule 2: \zh{不} before a 4th tone]
\zh{不} is normally \emph{bù} (4th tone). Before another 4th tone it becomes \emph{bú}.
\begin{itemize}
  \item \zh{不是} \emph{bú shì} ``is not''; \zh{不对} \emph{bú duì} ``not right'';
        \zh{不客气} \emph{bú kèqi}.
  \item but \zh{不好} \emph{bù hǎo}, \zh{不忙} \emph{bù máng}, \zh{不喝} \emph{bù hē}.
\end{itemize}
In the middle of a verb phrase (\zh{好不好} \emph{hǎo bu hǎo}) it is neutral.
\end{important}

\subsection{\zh{一} yī}

\begin{important}[Rule 3: \zh{一}]
\begin{itemize}
  \item Alone, when counting, in ordinals and in dates: \emph{yī} (\zh{第一} dì-yī,
        \zh{一月} yīyuè, \zh{十一} shíyī).
  \item Before a 4th tone (or a neutral tone that was originally 4th): \emph{yí}
        (\zh{一个} yí ge, \zh{一下} yíxià, \zh{一样} yíyàng).
  \item Before 1st, 2nd and 3rd tones: \emph{yì} (\zh{一天} yì tiān, \zh{一年} yì nián,
        \zh{一起} yìqǐ).
\end{itemize}
\end{important}

\begin{center}
\begin{tabular}{lllll}
\toprule
 & before 1st & before 2nd & before 3rd & before 4th \\
\midrule
\zh{不} & bù & bù & bù & \textbf{bú} \\
\zh{一} & \textbf{yì} & \textbf{yì} & \textbf{yì} & \textbf{yí} \\
3rd tone & half 3rd & half 3rd & \textbf{2nd} & half 3rd \\
\bottomrule
\end{tabular}
\end{center}

\begin{remark}
In our textbook \zh{不} and \zh{一} were printed with their \emph{changed} tones, but
the 3--3 sandhi was not marked. Quizzes accepted either for \zh{一} and \zh{不}, but
dictation required the dictionary tone for third tones.
\end{remark}

\section{Chinese characters}

A character (\zh{字} zì) stands for one syllable and (usually) one unit of meaning.
Modern words are mostly two characters: \zh{老师} ``teacher'' = \zh{老} ``old'' +
\zh{师} ``master''. An educated reader knows 3000--4000 characters; about 1000 cover
90\% of everyday text.

\subsection{How characters were formed}

Traditional scholarship classifies characters into six types (\zh{六书} liùshū); four of
them matter for learners.

\begin{center}
\begin{tabular}{llp{5.5cm}}
\toprule
Type & Examples & Idea \\
\midrule
Pictographs & \zh{人 山 日 月 木 口} & stylized drawings of the thing \\
Simple indicatives & \zh{一 二 三 上 下} & abstract signs (a mark above or below a line) \\
Compound ideographs & \zh{好 明 休 林} & two meanings combined: \zh{女} woman + \zh{子} child = good; \zh{日} sun + \zh{月} moon = bright; \zh{人} person + \zh{木} tree = rest \\
Phonetic--semantic & \zh{妈 吗 骂 们 问} & a \emph{radical} gives the meaning, another part gives the sound \\
\bottomrule
\end{tabular}
\end{center}

About 80\% of characters are phonetic--semantic compounds. The sound hint is often only
approximate (tones and initials drift), but it is still useful:

\begin{example}[The \zh{马} family]
\zh{马} mǎ ``horse'' is the sound component in
\zh{妈} mā (\zh{女} woman $\rightarrow$ mother),
\zh{吗} ma (\zh{口} mouth $\rightarrow$ spoken particle),
\zh{骂} mà (two mouths $\rightarrow$ scold),
\zh{码} mǎ (\zh{石} stone $\rightarrow$ number, code).
\end{example}

\section{Strokes}

Every character is written with a fixed set of strokes (\zh{笔画} bǐhuà).

\begin{center}
\begin{tabular}{llll}
\toprule
Stroke & Name & Direction & Appears in \\
\midrule
\zh{一} & \zh{横} héng & horizontal, left to right & \zh{三 王} \\
\zh{丨} & \zh{竖} shù & vertical, top to bottom & \zh{十 中} \\
\zh{丿} & \zh{撇} piě & falling to the left & \zh{人 八} \\
\zh{丶} & \zh{点} diǎn & dot & \zh{六 文} \\
\zh{㇏} & \zh{捺} nà & falling to the right & \zh{人 大} \\
\zh{㇀} & \zh{提} tí & rising to the right & \zh{我 打} \\
\zh{㇕} & \zh{横折} héngzhé & horizontal then turn down & \zh{口 日} \\
\zh{亅} & \zh{竖钩} shùgōu & vertical with hook & \zh{小 你} \\
\bottomrule
\end{tabular}
\end{center}

\subsection{Stroke order}

Stroke order matters for neatness, for counting strokes (dictionaries) and for
handwriting input on phones.

\begin{important}[Stroke order rules]
\begin{enumerate}
  \item Top before bottom: \zh{三}.
  \item Left before right: \zh{你}, \zh{明}.
  \item Horizontal before vertical when they cross: \zh{十}.
  \item Left-falling before right-falling: \zh{人}, \zh{八}.
  \item Outside before inside: \zh{月}, \zh{问}.
  \item Inside before closing the frame: \zh{国}, \zh{日} (the bottom stroke is last).
  \item Middle before the two sides: \zh{小}, \zh{水}.
  \item Enclosures open at the bottom-left (\zh{辶}, \zh{廴}) are written last:
        \zh{这}, \zh{还}.
\end{enumerate}
\end{important}

The diagrams show three basic characters stroke by stroke, with numbers at the start
of each stroke.

\begin{center}
\begin{tikzpicture}[stroke/.style={line width=2.2pt, line cap=round}, num/.style={circle, fill=red!75!black, text=white, inner sep=1pt, font=\scriptsize\bfseries}]
  % grid boxes
  \foreach \x/\lab in {0/{\zh{十} shí}, 3.5/{\zh{口} kǒu}, 7/{\zh{人} rén}} {
    \draw[gray] (\x,0) rectangle (\x+2.4,2.4);
    \draw[gray!40, dashed] (\x+1.2,0) -- (\x+1.2,2.4);
    \draw[gray!40, dashed] (\x,1.2) -- (\x+2.4,1.2);
    \node[below] at (\x+1.2,-0.05) {\lab};
  }
  % 十
  \draw[stroke] (0.3,1.25) -- (2.1,1.25);
  \draw[stroke] (1.2,2.15) -- (1.2,0.2);
  \node[num] at (0.15,1.55) {1};
  \node[num] at (0.9,2.2) {2};
  % 口
  \draw[stroke] (4.0,1.9) -- (4.0,0.5);
  \draw[stroke] (4.0,1.9) -- (5.4,1.9) -- (5.4,0.5);
  \draw[stroke] (4.0,0.55) -- (5.4,0.55);
  \node[num] at (3.75,1.55) {1};
  \node[num] at (4.7,2.15) {2};
  \node[num] at (4.7,0.25) {3};
  % 人
  \draw[stroke] (8.3,2.1) .. controls (8.2,1.2) .. (7.4,0.3);
  \draw[stroke] (8.25,1.4) .. controls (8.6,0.8) .. (9.1,0.3);
  \node[num] at (8.6,2.15) {1};
  \node[num] at (8.75,1.35) {2};
\end{tikzpicture}
\end{center}

\section{Radicals}

The \emph{radical} (\zh{部首} bùshǒu) is the component under which a character is
listed in a dictionary, and usually hints at meaning. Many radicals change shape when
they are squeezed to the side of a character.

\begin{center}
\begin{tabular}{lllll}
\toprule
Radical & Side form & Name / meaning & Examples \\
\midrule
\zh{人} & \zh{亻} & person & \zh{你 他 们 做} \\
\zh{女} & \zh{女} & woman & \zh{妈 姐 好 她} \\
\zh{口} & \zh{口} & mouth & \zh{吗 叫 喝 吃} \\
\zh{水} & \zh{氵} & water & \zh{汉 没 洗 酒} \\
\zh{言} & \zh{讠} & speech & \zh{说 话 请 谢 语} \\
\zh{手} & \zh{扌} & hand & \zh{打 找 拿} \\
\zh{心} & \zh{忄} & heart & \zh{忙 快 您 想} \\
\zh{木} & \zh{木} & tree, wood & \zh{校 机 椅} \\
\zh{日} & \zh{日} & sun, day & \zh{明 时 早 星} \\
\zh{丝} & \zh{纟} & silk, thread & \zh{红 绿 给 经} \\
\zh{草} & \zh{艹} & grass & \zh{茶 菜 英 花} \\
\zh{走} & \zh{辶} & walking & \zh{这 还 进 远} \\
\zh{门} & \zh{门} & door & \zh{们 问 间} \\
\bottomrule
\end{tabular}
\end{center}

\section{Character structure}

Each character fits in an imaginary square. Its components are arranged in a small
number of patterns; recognising the pattern makes new characters much easier to
memorize and write.

\begin{center}
\begin{tikzpicture}[comp/.style={draw, thick, fill=blue!8, align=center, font=\large}]
  % left-right 好
  \draw[thick] (0,0) rectangle (2,2);
  \node[comp, minimum width=0.9cm, minimum height=1.8cm] at (0.5,1) {\zh{女}};
  \node[comp, fill=orange!15, minimum width=0.9cm, minimum height=1.8cm] at (1.5,1) {\zh{子}};
  \node[below, align=center, font=\small] at (1,-0.1) {left--right\\ \zh{好}};
  % top-bottom 字
  \draw[thick] (2.8,0) rectangle (4.8,2);
  \node[comp, inner sep=0pt, font=\normalsize, minimum width=1.8cm, minimum height=0.55cm] at (3.8,1.6) {\zh{宀}};
  \node[comp, fill=orange!15, minimum width=1.8cm, minimum height=1.0cm] at (3.8,0.62) {\zh{子}};
  \node[below, align=center, font=\small] at (3.8,-0.1) {top--bottom\\ \zh{字}};
  % full enclosure 国
  \draw[thick] (5.6,0) rectangle (7.6,2);
  \draw[line width=3pt, blue!60] (5.75,0.15) rectangle (7.45,1.85);
  \node[comp, fill=orange!15, minimum width=0.9cm, minimum height=0.9cm] at (6.6,1) {\zh{玉}};
  \node[below, align=center, font=\small] at (6.6,-0.1) {enclosure\\ \zh{国} = \zh{囗} + \zh{玉}};
  % half enclosure 这
  \draw[thick] (8.4,0) rectangle (10.4,2);
  \draw[line width=3pt, blue!60] (8.6,1.9) -- (8.6,0.2) -- (10.3,0.2);
  \node[comp, fill=orange!15, minimum width=1.2cm, minimum height=1.2cm] at (9.6,1.15) {\zh{文}};
  \node[below, align=center, font=\small] at (9.4,-0.1) {lower-left\\ \zh{这} = \zh{辶} + \zh{文}};
  % left-middle-right 谢
  \draw[thick] (11.2,0) rectangle (13.2,2);
  \node[comp, minimum width=0.5cm, minimum height=1.8cm] at (11.5,1) {\zh{讠}};
  \node[comp, fill=orange!15, minimum width=0.6cm, minimum height=1.8cm] at (12.15,1) {\zh{身}};
  \node[comp, fill=green!15, minimum width=0.5cm, minimum height=1.8cm] at (12.85,1) {\zh{寸}};
  \node[below, align=center, font=\small] at (12.2,-0.1) {left--middle--right\\ \zh{谢}};
\end{tikzpicture}
\end{center}

\section{Simplified and traditional characters}

The PRC simplified about 2000 characters in the 1950s and 1960s; Singapore followed.
Taiwan, Hong Kong and many older overseas communities use traditional characters. This
course uses simplified characters, but we should be able to recognise the traditional
forms of very common ones.

\begin{center}
\begin{tabular}{llll}
\toprule
Simplified & Traditional & Pinyin & Meaning \\
\midrule
\zh{学} & \zht{學} & xué & to study \\
\zh{国} & \zht{國} & guó & country \\
\zh{汉} & \zht{漢} & hàn & Han (Chinese) \\
\zh{语} & \zht{語} & yǔ & language \\
\zh{谢} & \zht{謝} & xiè & to thank \\
\zh{们} & \zht{們} & men & plural suffix \\
\zh{马} & \zht{馬} & mǎ & horse \\
\zh{门} & \zht{門} & mén & door \\
\bottomrule
\end{tabular}
\end{center}

\begin{aside}
Simplification often replaced a whole component with a shorthand that was already used
in handwriting: \zht{言} $\rightarrow$ \zh{讠}, \zht{門} $\rightarrow$ \zh{门},
\zht{馬} $\rightarrow$ \zh{马}. Learning the radical simplifications lets you read
dozens of traditional characters at once. Many characters were never simplified at all
(\zh{人 山 好 我 你}).
\end{aside}

\section*{Review}
\begin{itemize}
  \item 3 + 3 $\rightarrow$ 2 + 3; a third tone before other tones is a low ``half third''.
  \item \zh{不} becomes \emph{bú} before a 4th tone; \zh{一} becomes \emph{yí} before a 4th
        tone and \emph{yì} before the others, but stays \emph{yī} when counting.
  \item Most characters are phonetic--semantic compounds: radical for meaning, another
        component for sound.
  \item Learn the basic strokes and the stroke-order rules (top--bottom, left--right,
        horizontal before vertical, outside before inside, close the frame last).
  \item Main structures: left--right, top--bottom, enclosure, half enclosure.
  \item Recognise traditional forms of very common characters.
\end{itemize}
